%% kapa: exam=1 round=26 year=2015 period=1 form=A subject=civil_law subject_ko=민법 seq=19 number=19
%% source: Q-Net 감정평가사 기출문제 (cst00309) artlSeq=5138001 file=extracted/1008/5138001/제26회_감정평가사 1차_기출/제26회_감정평가사 1차_기출/1교시_감평_A형.pdf
%% conversion: kordoc 4.24.51 -> tools/gen1.py (auto); review=proofread-sonnet; flags=none
\begin{question}{19}
소멸시효와 제척기간에 관한 설명으로 옳은 것은? (다툼이 있으면 판례에 따름)
\begin{choices}
  \item 소멸시효와 제척기간 모두 중단과 정지가 인정된다.
  \item 소멸시효의 중단에 관한 규정은 취득시효에 준용한다.
  \item 소멸시효는 법률행위로 단축할 수 없다.
  \item 당사자가 본래의 소멸시효 기산일보다 뒤의 날짜를 기산일로 주장하는 경우 법원은 본래의 소멸시효 기산일을 기준으로 소멸시효를 계산하여야 한다.
  \item 제척기간은 소송상 당사자가 제척기간의 도과를 주장한 경우에 한하여 고려된다.
\end{choices}
\end{question}
\answer{2}
